% TOP_PROBLEM: {"id": 1398, "title": "Jaeger's Petersen Coloring Conjecture", "category_id": 3, "category": {"id": 3, "name": "graph_theory", "display_name": "Graph Theory", "description": "Problems involving graphs, networks, and their properties.", "slug": "graph-theory", "order_index": 3, "created_at": "2026-07-31T15:26:25.671Z"}, "status": "open"}
% FIRST_POSTED: null
% ENTRY_KIND: statement_only
% Imported from: batch06.tex; UnsolvedMath ID 1398
\documentclass[11pt]{article}
\usepackage[margin=25mm]{geometry}
\usepackage{fontspec}
\setmainfont{Latin Modern Roman}
\usepackage{amsmath,amssymb,mathtools}
\usepackage{unicode-math}
\setmathfont{Latin Modern Math}
\usepackage{xurl,hyperref}
\hypersetup{colorlinks=true,urlcolor=blue}
\setcounter{secnumdepth}{0}
\setlength{\parindent}{0pt}
\setlength{\parskip}{5pt}
\setlength{\emergencystretch}{3em}
\newcommand{\sref}[2]{\hyperlink{source:#1:#2}{[S#2]}}
\newcommand{\cataloguescope}[1]{\paragraph{Recorded target.}#1}
\begin{document}
% UnsolvedMath ID 1398; source catalogue code GRAPH-011
\setcounter{section}{1397}
\section{Jaeger's Petersen Coloring Conjecture}\label{1398}
{\small\sffamily UnsolvedMath ID 1398 · Graph Theory}
\subsection{Definitions and mathematical statement}

\paragraph{Shared notation.}$\mathbb N=\{1,2,\ldots\}$, and $\mathbb Z,\mathbb Q,\mathbb R,\mathbb C$ denote the integers, rationals, reals and complex numbers. Any other conventions are specified locally.


A finite loopless undirected graph is cubic if every vertex has degree three, and bridgeless if deleting any edge does not increase its number of components. The Petersen graph $P$ is the Kneser graph on the two-element subsets of $\{1,2,3,4,5\}$, with disjoint subsets adjacent. For a vertex $v$, $\delta_G(v)$ is the set of incident edges. A Petersen coloring is a map $f:E(G)\to E(P)$ for which each incident triple $\delta_G(v)$ maps bijectively to $\delta_P(w)$ for some $w\in V(P)$. Such an edge map is cycle-continuous: the inverse image of the edge set of any cycle of $P$ has even degree at every vertex of $G$.
\paragraph{Statement recorded in the supplied source.}
Does every finite bridgeless cubic graph $G$ admit a Petersen coloring $f:E(G)\to E(P)$? In the equivalent normal-coloring formulation, does $G$ have a proper edge coloring $c:E(G)\to\{1,2,3,4,5\}$ such that for each edge $e=uv$,
\[
 \left|\{c(a):a\in(\delta_G(u)\cup\delta_G(v))\setminus\{e\}\}\right|\in\{2,4\}?
\]
\sref{1398}{2}
\subsection{Short English statement}
Can every bridgeless cubic graph map its incident edge triples to incident triples of the Petersen graph, equivalently admit a normal proper edge coloring with five colors?
\subsection{Sources}
\begin{description}
\item[\hypertarget{source:1398:1}{S1}] ULAM.AI and the UnsolvedMath Contributors, \emph{UnsolvedMath: A Curated Collection of Open Mathematics Problems}, record 1398 (GRAPH-011). CC BY 4.0. \url{https://huggingface.co/datasets/ulamai/UnsolvedMath}.
\item[\hypertarget{source:1398:2}{S2}] François Pirot, Jean-Sébastien Sereni and Riste Škrekovski, \emph{Variations on the Petersen colouring conjecture} (2019), arXiv:1905.07913v1, Section 1, Conjecture 1.1 and Proposition 1.2. \url{https://arxiv.org/abs/1905.07913v1}.
\end{description}
\label{end:1398}
\end{document}
